\documentclass[11pt,a4paper]{article}

\usepackage[T1]{fontenc}
\usepackage[utf8]{inputenc}
% babel dutch omitted for portable MiKTeX builds
\usepackage[margin=2.2cm]{geometry}
\usepackage{amsmath,amssymb}
\usepackage{booktabs}
\usepackage{graphicx}
\usepackage{hyperref}
\usepackage{enumitem}
\usepackage{array}
\usepackage{multirow}

\title{Toy walkthrough: spectrale SIS-drempel op een mini-netwerk\\[0.4em]
\large Handmatige uitwerking bij het grp-SIS-onderzoeksproject}
\author{Jan Ruijgrok \textbar\ Open Universiteit \textbar\ IM1312}
\date{Juli 2026}

\begin{document}
\maketitle

\begin{abstract}
Dit document is een \textbf{didactische versie} van het onderzoeksproject \emph{Testing a spectral epidemic threshold for SIS on networks under heavy-tailed recovery}. Alle stappen worden doorlopen met \textbf{kleine getallen} die je met pen en papier kunt nakijken. De logica is identiek aan het echte project (Erd\H{o}s--R\'enyi, $N{\approx}2000$, BehaviorSpace), maar het netwerk heeft slechts 6 knopen en 6 replicates per infectiekans~$\beta$.
\end{abstract}

\section*{Onderzoeksvraag in \'e\'en zin}
\begin{quote}
Klopt de spectrale drempel $\tau_c \approx 1/\lambda_{\max}$ (en de bijbehorende $\beta_{\mathrm{pred}}$) nog als we recovery-tijden \textbf{zwaar-staartig} maken, terwijl het \textbf{gemiddelde} $\mathbb{E}[W]$ gelijk blijft?
\end{quote}

\section{Overzicht: welke stappen doorrekenen?}

\begin{center}
\small
\begin{tabular}{@{}clll@{}}
\toprule
\textbf{Stap} & \textbf{Wat} & \textbf{Toy (dit document)} & \textbf{Echt project} \\
\midrule
1 & Netwerk opbouwen & 6-cyclus (ring) & ER, $N{=}2000$, $\langle k\rangle{\approx}6$ \\
2 & $\lambda_{\max}$ & cosinusformule & Python op edge-CSV \\
3 & Spectrale voorspelling & $\beta_{\mathrm{pred}}=1/(\lambda_{\max}\mathbb{E}[W])$ & idem, $\lambda_{\max}{\approx}7{,}18$ \\
4 & Hersteltijd $\mathbb{E}[W]{=}5$ & exp.\ / power-law / lognormal & zelfde drie regimes \\
5 & Simulaties per $\beta$ & vaste toy-tabel (6 reps) & BehaviorSpace, 24 reps \\
6--7 & Drempel $\hat\beta_{\mathrm{surv}\,50}$ & survival $\ge 50\%$ & idem + bootstrap CI \\
8 & Extinctietijden (RQ3) & mediaan tick $|$ extinct & idem in rapport \\
\bottomrule
\end{tabular}
\end{center}

\section{Stap 1 --- Klein contactnetwerk}

In het SIS-model (susceptible--infected--susceptible) kan een agent na herstel \textbf{opnieuw} besmet raken. Infectie verspreidt zich langs \textbf{contacten} (randen). De vraag is: bij welke infectiekracht $\beta$ sterft een uitbraak uit, en wanneer blijft ze hangen?

We gebruiken een \textbf{6-cyclus} (ring): elke knoop heeft precies 2 buren.

\begin{center}
\texttt{0 --- 1 --- 2}\\
\texttt{| \hspace{3.2em} |}\\
\texttt{5 --- 4 --- 3}
\end{center}

\textbf{Randen (6 stuks):} $(0,1)$, $(1,2)$, $(2,3)$, $(3,4)$, $(4,5)$, $(5,0)$.

\subsection*{Vaste parameters}

\begin{center}
\begin{tabular}{@{}llll@{}}
\toprule
Parameter & Betekenis & Toy & Echt project \\
\midrule
$N$ & aantal agents & 6 & 2000 \\
$\mathbb{E}[W]$ & gem.\ hersteltijd (ticks) & 5 & 5 \\
$I_0$ & initieel besmet & 1 & 5 \\
max-ticks & simulatie-horizon & 500 & 10\,000 \\
$\beta$ & infectiekans per tick langs link & wordt gevarieerd & sweep 0.018--0.056 \\
\bottomrule
\end{tabular}
\end{center}

\subsection*{Adjacency-matrix $A$}

\begin{center}
\[
A =
\begin{pmatrix}
0 & 1 & 0 & 0 & 0 & 1 \\
1 & 0 & 1 & 0 & 0 & 0 \\
0 & 1 & 0 & 1 & 0 & 0 \\
0 & 0 & 1 & 0 & 1 & 0 \\
0 & 0 & 0 & 1 & 0 & 1 \\
1 & 0 & 0 & 0 & 1 & 0
\end{pmatrix}
\]
\end{center}

\textbf{Controle:} 6 randen, 6 knopen $\Rightarrow$ $\langle k\rangle = 2m/N = 12/6 = \mathbf{2}$.

\section{Stap 2 --- $\lambda_{\max}$ uit de adjacency-matrix}

De \textbf{grootste eigenwaarde} $\lambda_{\max}$ van $A$ hangt samen met hoe gemakkelijk infectie door het netwerk kan blijven circuleren. Voor een $N$-cyclus (circulant matrix):

\[
\lambda_j = 2\cos\!\left(\frac{2\pi j}{N}\right), \quad j=0,1,\ldots,N-1.
\]

Voor $N=6$:

\begin{center}
\begin{tabular}{@{}cc@{}}
\toprule
$j$ & $\lambda_j$ \\
\midrule
0 & $2\cos(0^\circ) = \mathbf{2}$ \\
1 & $2\cos(60^\circ) = 1$ \\
2 & $2\cos(120^\circ) = -1$ \\
3 & $-2$ \\
4 & $-1$ \\
5 & $1$ \\
\bottomrule
\end{tabular}
\end{center}

Dus \textbf{$\lambda_{\max} = 2$}.

\subsection*{Graden (controle + mean-field alternatief)}

\begin{itemize}[nosep]
  \item $\langle k\rangle = 2$
  \item $\langle k^2\rangle = 4$ (iedereen graad~2)
  \item Homogene MF: $\tau_{\mathrm{hom}} = 1/\langle k\rangle = 0{,}50$
  \item Heterogene MF: $\tau_{\mathrm{het}} = \langle k\rangle/\langle k^2\rangle = 0{,}50$
\end{itemize}

Op deze reguliere ring vallen spectrale en MF-referenties samen. Op jullie ER-graaf schelen ze iets ($\beta_{\mathrm{pred}}{\approx}0{,}028$ vs.\ $\beta_{\mathrm{het}}{\approx}0{,}029$).

\section{Stap 3 --- Spectrale voorspelling $\beta_{\mathrm{pred}}$}

Effectieve transmissie-intensiteit:
\[
\tau = \beta \cdot \mathbb{E}[W].
\]

Mean-field spectrale referentie ($c=1$):
\[
\tau_{\mathrm{pred}} = \frac{1}{\lambda_{\max}} = \frac{1}{2} = 0{,}50,
\qquad
\beta_{\mathrm{pred}} = \frac{\tau_{\mathrm{pred}}}{\mathbb{E}[W]} = \frac{0{,}50}{5} = \mathbf{0{,}10}.
\]

\textbf{Interpretatie:} onder de mean-field benchmark zou de grens tussen uitsterven en persisteren \emph{rond} $\beta \approx 0{,}10$ liggen. In het echte project test je of de empirische drempel $\hat\beta_{\mathrm{surv}\,50}$ daar dichtbij ligt (RQ1, exponentieel) of verschuift (RQ2, zware staarten).

\section{Stap 4 --- Hersteltijdverdelingen (grp-SIS)}

Het grp-SIS-kader (Tang et al.) laat de \textbf{vorm} van de hersteltijd $W$ vrij, maar houdt $\mathbb{E}[W]=5$ vast in alle drie regimes:

\begin{center}
\begin{tabular}{@{}lll@{}}
\toprule
Regime & Karakter & Illustratieve draws $W$ (ticks) \\
\midrule
Exponentieel & geheugenloos & 3, 7, 4, 6, \ldots \\
Power-law (Tang) & zware rechterstaart & 2, 3, \textbf{45}, 4, \ldots \\
Lognormal (Tang) & brede spreiding & 1, 8, \textbf{30}, 5, \ldots \\
\bottomrule
\end{tabular}
\end{center}

\textbf{Belangrijk:} omdat $\tau = \beta \cdot \mathbb{E}[W]$ en $\mathbb{E}[W]$ vast is, verandert $\beta_{\mathrm{pred}}$ \emph{niet} tussen regimes. Eventuele verschuivingen in $\hat\beta_{\mathrm{surv}\,50}$ komen door de \textbf{dynamiek}, niet doordat het gemiddelde herstel verschuift.

\section{Stap 5 --- Simulatie-uitkomsten (toy BehaviorSpace)}

Per run noteer je:
\begin{itemize}[nosep]
  \item \texttt{bs-out-extinct} $= 1$ als alles genezen is v\'o\'or max-ticks; anders $0$
  \item \texttt{bs-out-final-tick} als extinct $= 1$
\end{itemize}

We gebruiken \textbf{6 replicates} per $(\beta, \text{regime})$. Hieronder: 1 = overleefd, 0 = uitgestorven.

\subsection*{Exponentieel herstel}

\begin{center}
\begin{tabular}{@{}ccc@{}}
\toprule
$\beta$ & Overleefd? (6 runs) & Survival rate \\
\midrule
0,06 & 0 0 0 0 0 0 & 0/6 = 0\% \\
0,08 & 0 0 0 0 0 1 & 1/6 = 17\% \\
0,10 & 0 0 0 1 0 1 & 2/6 = 33\% \\
0,12 & 1 1 0 1 1 0 & 4/6 = \textbf{67\%} \\
0,14 & 1 1 1 1 1 1 & 6/6 = 100\% \\
\bottomrule
\end{tabular}
\end{center}

\subsection*{Power-law herstel}

\begin{center}
\begin{tabular}{@{}ccc@{}}
\toprule
$\beta$ & Overleefd? & Survival rate \\
\midrule
0,06 & 0 0 0 0 0 0 & 0\% \\
0,08 & 0 1 0 1 0 0 & 33\% \\
0,10 & 1 0 1 1 0 0 & \textbf{50\%} \\
0,12 & 1 1 1 1 0 1 & 83\% \\
0,14 & 1 1 1 1 1 1 & 100\% \\
\bottomrule
\end{tabular}
\end{center}

\subsection*{Lognormal herstel}

\begin{center}
\begin{tabular}{@{}ccc@{}}
\toprule
$\beta$ & Overleefd? & Survival rate \\
\midrule
0,10 & 0 0 0 0 0 1 & 17\% \\
0,12 & 0 0 1 0 1 0 & 33\% \\
0,14 & 1 1 0 1 1 0 & \textbf{67\%} \\
\bottomrule
\end{tabular}
\end{center}

\section{Stap 6 \& 7 --- Survival curve en drempel $\hat\beta_{\mathrm{surv}\,50}$}

\textbf{Regel} (zoals in \texttt{scripts/threshold\_estimators.py}):

\begin{quote}
$\hat\beta_{\mathrm{surv}\,50}$ = \textbf{kleinste} $\beta$ op het grid waarvoor survival rate $\ge 50\%$.
\end{quote}

\begin{center}
\begin{tabular}{@{}lccc@{}}
\toprule
Regime & Eerste $\beta$ met $\ge 50\%$ & $\hat\beta_{\mathrm{surv}\,50}$ & $\hat\beta / \beta_{\mathrm{pred}}$ \\
\midrule
Exponentieel & 0,12 (67\%) & \textbf{0,12} & 1,20 \\
Power-law & 0,10 (50\%) & \textbf{0,10} & 1,00 \\
Lognormal & 0,14 (67\%) & \textbf{0,14} & 1,40 \\
\bottomrule
\end{tabular}
\end{center}

In $\tau$-taal ($\hat\tau_{50} = \hat\beta \cdot 5$, $\tau_{\mathrm{pred}} = 0{,}50$):

\begin{center}
\begin{tabular}{@{}ccc@{}}
\toprule
Regime & $\hat\tau_{50}$ & $\hat\tau_{50}/\tau_{\mathrm{pred}}$ \\
\midrule
Exponentieel & 0,60 & 1,20 \\
Power-law & 0,50 & 1,00 \\
Lognormal & 0,70 & 1,40 \\
\bottomrule
\end{tabular}
\end{center}

\subsection*{Koppeling aan onderzoeksvragen}

\begin{itemize}[leftmargin=*]
  \item \textbf{RQ1} (exponentieel): empirische drempel (0,12) ligt \textbf{boven} spectrale voorspelling (0,10) $\Rightarrow$ ratio 1,2. In het echte project: ratio $\approx 1{,}65$.
  \item \textbf{RQ2}: power-law drempel \textbf{lager} (0,10), lognormal \textbf{hoger} (0,14) dan exponentieel (0,12).
  \item \textbf{RQ3}: drempels verschillen matig; extinctietijd onderscheidt regimes sterker (stap~8).
\end{itemize}

\section{Stap 8 --- Extinctietijden (RQ3)}

Bij $\beta = 0{,}12$, alleen \textbf{uitgestorven} runs:

\begin{center}
\begin{tabular}{@{}lccc@{}}
\toprule
Regime & Uitgestorven runs & Extinctieticks & Mediaan \\
\midrule
Exponentieel & 2 & 52, 45 & \textbf{48,5} \\
Power-law & 1 & 28 & \textbf{28} \\
Lognormal & 4 & 30, 25, 28, 22 & \textbf{27} \\
\bottomrule
\end{tabular}
\end{center}

\textbf{Interpretatie:} als een uitbraak uitsterft, gebeurt dat onder zware-staart-herstel \textbf{sneller} dan onder exponentieel --- terwijl power-law bij \emph{lagere} $\beta$ juist \textbf{vaker} overleeft. Drempel en extinctietijd zijn dus \textbf{verschillende} meetinstrumenten (vergelijkbaar met het rapport: median extinct tick 22 vs.\ 17 vs.\ 15).

\section{Vergelijking met het echte project}

\begin{center}
\begin{tabular}{@{}lcc@{}}
\toprule
 & Toy ring & Echt ER (seed 10001) \\
\midrule
$N$ & 6 & 2000 \\
$\lambda_{\max}$ & 2 & 7,18 \\
$\beta_{\mathrm{pred}}$ & 0,10 & 0,028 \\
$\hat\beta$ exponentieel & 0,12 (ratio 1,20) & 0,046 (ratio 1,65) \\
$\hat\beta$ power-law & 0,10 (ratio 1,00) & 0,044 (ratio 1,58) \\
$\hat\beta$ lognormal & 0,14 (ratio 1,40) & 0,050 (ratio 1,80) \\
\bottomrule
\end{tabular}
\end{center}

\textbf{Wat blijft gelijk:} empirische drempel boven spectrale lijn; power-law laagste, lognormal hoogste; RQ3-signalen in extinctietijden.

\textbf{Wat verschilt:} absolute ratios door $N$, discrete tijd en gridfijnheid ($\Delta\beta = 0{,}02$ hier vs.\ 0,002 in het rapport).

\section{Inference-niveaus (Wieringa / IM1312)}

\begin{center}
\small
\begin{tabular}{@{}p{2.2cm}p{4.2cm}p{7.5cm}@{}}
\toprule
\textbf{Niveau} & \textbf{Vraag} & \textbf{Voorbeeld} \\
\midrule
(a) Descriptief & Wat gebeurde in onze runs? & ``Bij $\beta{=}0{,}12$ overleefden 4/6 exponenti\"ele runs.'' \\
(d) Statistisch & Generaliseert dit? & Bootstrap CI; meerdere ER-seeds in rapport \\
(b) Abductief & Plausibele verklaring? & ``Gap door discrete tijd + eindige $N$.'' \\
(c) Analogisch & Geldt dit elders? & ``Verwacht shift op lattice/ring'' (nog testen) \\
\bottomrule
\end{tabular}
\end{center}

\vspace{0.5em}
\textbf{Fout:} \emph{``Power-law recovery veroorzaakt 20\% lagere drempels.''} (causaal, niet bewezen)

\textbf{Goed:} \emph{``Power-law $\hat\beta$ lag 0,02 onder exponentieel op dezelfde toy-graaf; consistent met langere infectie-staarten, maar $N{=}6$ en grof grid beperken de precisie.''}

\section*{Oefeningen en antwoordsleutel}

\subsection*{Opdrachten (pen en papier)}
\begin{enumerate}[leftmargin=*]
  \item Bereken alle zes eigenwaarden van de 6-cyclus en bevestig $\lambda_{\max}=2$.
  \item Bereken $\beta_{\mathrm{pred}}$ met $\mathbb{E}[W]=5$.
  \item Exponentieel, $\beta=0{,}12$: hoeveel runs hebben \texttt{extinct}$=0$? Wat is de survival rate?
  \item Bepaal $\hat\beta_{\mathrm{surv}\,50}$ voor elk van de drie regimes.
  \item Bereken $\hat\tau_{50} = \hat\beta_{\mathrm{surv}\,50} \cdot 5$ en $\hat\tau_{50}/\tau_{\mathrm{pred}}$ voor exponentieel.
  \item Bij $\beta=0{,}12$: mediaan extinctietick voor exponentieel (alleen uitgestorven runs).
  \item Welke conclusies blijven gelijk in het echte ER-project?
\end{enumerate}

\subsection*{Antwoordsleutel}
\begin{center}
\begin{tabular}{@{}cl@{}}
\toprule
\# & Antwoord \\
\midrule
1 & $\lambda \in \{2, 1, -1, -2, -1, 1\}$ $\Rightarrow$ max $= 2$ \\
2 & $\beta_{\mathrm{pred}} = 1/(2 \times 5) = 0{,}10$ \\
3 & 4 overleefd $\Rightarrow$ survival $= 4/6 \approx$ \textbf{67\%} \\
4 & exp.\ \textbf{0,12}; power-law \textbf{0,10}; lognormal \textbf{0,14} \\
5 & $\hat\tau_{50} = 0{,}60$; ratio $0{,}60/0{,}50 = \mathbf{1{,}20}$ \\
6 & ticks 52 en 45 $\Rightarrow$ mediaan $(52+45)/2 =$ \textbf{48,5} \\
7 & Drempel boven spectrale lijn; power-law laagste, lognormal hoogste; extinctietijden onderscheiden regimes sterker \\
\bottomrule
\end{tabular}
\end{center}

\vspace{1em}
\hrule
\vspace{0.5em}
\small
\textit{Bijbehorend Jupyter-notebook: \texttt{notebooks/toy\_sis\_threshold\_walkthrough.ipynb}. Projectrepo: master-rmai-researchproject-virussimulation.}

\end{document}
